% Bagian: first-order-logic
% Bab: beyond
% Subbagian: modal-logics

\documentclass[../../../include/open-logic-section]{subfiles}

\begin{document}

\olfileid[id]{fol}{byd}{mod}

\olsection{Logika Modal}

Pertimbangkan contoh kalimat bersyarat berikut:
\begin{quote}
  Jika Jeremy sendirian di ruangan itu, maka ia mabuk dan telanjang dan
  menari-nari di atas kursi.
\end{quote}
Ini merupakan contoh pernyataan bersyarat yang mungkin benar secara material,
tetapi tetap menyesatkan, karena tampaknya menyiratkan bahwa terdapat hubungan
yang lebih kuat antara anteseden dan kesimpulannya daripada sekadar bahwa
anteseden itu salah atau konsekuennya benar. Artinya, susunan katanya
menyiratkan bahwa klaim tersebut bukan hanya benar di dunia khusus ini (tempat
klaim itu mungkin benar secara trivial karena Jeremy tidak sendirian di
ruangan tersebut), tetapi lebih jauh bahwa kesimpulannya \emph{akan} benar
\emph{seandainya} antesedennya benar. Dengan kata lain, pernyataan itu dapat
dipahami sebagai mengatakan bahwa klaim tersebut benar bukan hanya di dunia
ini, melainkan di setiap dunia yang ``mungkin''; atau bahwa klaim itu
\emph{niscaya} benar, bukan sekadar benar di dunia khusus ini.

Logika modal dirancang untuk menjelaskan jenis keniscayaan ini. Logika
proposisional modal diperoleh dari logika proposisional biasa dengan
menambahkan operator kotak; yakni, jika $!A$ merupakan !!a{formula}, demikian
pula $\Box !A$. Secara intuitif, $\Box !A$ menyatakan
bahwa $!A$ \emph{niscaya} benar, atau benar di setiap dunia mungkin.
$\Diamond !A$ biasanya dianggap sebagai singkatan bagi $\lnot \Box \lnot
!A$, dan dapat dibaca sebagai menyatakan bahwa $!A$ \emph{mungkin} benar.
Tentu saja, modalitas juga dapat ditambahkan pada logika predikat.

Berbagai !!{structure}s Kripke dapat digunakan untuk menyediakan semantik bagi logika
modal; bahkan, Kripke pertama kali merancang semantik ini dengan logika modal
dalam pikirannya. Alih-alih membatasi diri pada urutan parsial, secara lebih
umum kita mempunyai himpunan ``dunia mungkin''~$P$ dan relasi
``aksesibilitas'' biner $\Atom{R}{x,y}$ di antara dunia-dunia itu. Secara
intuitif, $\Atom{R}{p,q}$ menyatakan bahwa dunia~$q$ selaras dengan~$p$;
yakni, jika kita ``berada di'' dunia~$p$, kita harus mempertimbangkan
kemungkinan bahwa dunia dapat saja seperti~$q$.

Logika modal kadang-kadang disebut logika ``intensional'', sebagai lawan dari
logika ``ekstensional''. Semantik yang dimaksudkan bagi logika ekstensional,
seperti logika klasik, hanya merujuk pada satu dunia, yakni dunia ``aktual'';
sementara semantik bagi logika ``intensional'' bertumpu pada ontologi yang
lebih rumit. Selain memodelkan keniscayaan, kita dapat menggunakan
modalitas untuk memodelkan konstruksi linguistik lain,
dengan menafsirkan ulang $\Box$ dan $\Diamond$ sesuai penerapannya. Misalnya:
\begin{enumerate}
\item Dalam logika keterbuktian, $\Box !A$ dibaca ``$!A$ dapat dibuktikan''
  dan $\Diamond !A$ dibaca ``$!A$ konsisten.''
\item Dalam logika epistemik, $\Box !A$ dapat dibaca sebagai ``Saya mengetahui
  $!A$'' atau ``Saya meyakini $!A$.''
\item Dalam logika temporal, $\Box !A$ dapat dibaca sebagai ``$!A$ selalu
  benar'' dan $\Diamond !A$ sebagai ``$!A$ kadang-kadang benar.''
\end{enumerate}

Kita ingin memperluas logika dengan aturan dan aksioma yang menangani
modalitas. Misalnya, sistem $\Log{S4}$ terdiri atas aksioma dan aturan logika
proposisional biasa, beserta aksioma-aksioma berikut:
\begin{align*}
& \Box (!A \lif !B) \lif (\Box !A \lif \Box !B)\\
& \Box !A \lif !A \\
& \Box !A \lif \Box \Box !A
\intertext{serta suatu aturan, ``dari $!A$, simpulkan $\Box !A$.''
  $\Log{S5}$ menambahkan aksioma berikut:}
& \Diamond !A \lif \Box \Diamond !A
\end{align*}
Variasi aksioma-aksioma ini mungkin cocok bagi penerapan yang berbeda-beda;
misalnya, S5 biasanya dianggap mencirikan gagasan keniscayaan logis. Hal yang
menarik ialah bahwa biasanya kita dapat menemukan semantik yang menjadikan
sistem !!{derivation} sahih dan lengkap dengan membatasi relasi aksesibilitas
dalam !!{structure}s Kripke dengan cara yang alami. Misalnya, $\Log{S4}$
bersesuaian dengan kelas !!{structure}s Kripke yang relasi aksesibilitasnya
refleksif dan transitif. $\Log{S5}$ bersesuaian dengan kelas !!{structure}s
Kripke yang relasi aksesibilitasnya {\em universal}, yakni setiap dunia dapat
diakses dari setiap dunia lain; sehingga $\Box !A$ berlaku jika dan hanya jika
$!A$ berlaku di setiap dunia.

\end{document}
